% !TeX root = ../../Book3.tex
% 纲要标题：### 4.1 切向与法向分解
% 纲要位置：几何学计划纲要.md，第 2230 行。
% 正文状态：初稿；尚未正式审核。

% 纲要内容开始
% 
% * 第二基本形式和形算子
% * 法联络及其曲率
% * 平均曲率向量
% * 全测地和全脐子流形
% 
% 纲要内容结束

\section{切向与法向分解}
\label{g3:sec:0401}

沿子流形移动一个切向量，环境中的导数可能离开切空间。第一章将导数的切向部分识别为诱导度量的 Levi-Civita 联络，现在保留下来的法向部分记录浸入怎样弯曲。法向量自身也会变化：它的切向导数由同一弯曲数据决定，法向导数则构成法丛上的联络。余维大于一时，法向空间中还有独立的转动，必须把这部分一起记录。

\subsection{第二基本形式和形算子}
\label{g3:subsec:0401:01}

设 \(f:(M^n,g)\to(\bar M^{n+r},\bar g)\) 为光滑等距浸入，其中 \(n\ge1\)，两流形均无边界，且 \(g=f^*\bar g\)。沿 \(f\) 的向量场，是给每个源点 \(p\) 指定一个向量 \(W(p)\in\mathrm{T}_{f(p)}\bar M\) 的光滑截面。它们所在的丛记为 \(f^*\mathrm{T}\bar M\)。即使两个不同源点有同一个像，这两个源点上的纤维仍分别记录。

\begin{definition}\label{g3:def:normal-bundle}
	设 \(f:(M^n,g)\to(\bar M^{n+r},\bar g)\) 为光滑等距浸入。其\term[facong]{法丛}[normal bundle]为
	\[
		N_fM\coloneqq
		\{(p,\xi):p\in M,\ \xi\in\mathrm{T}_{f(p)}\bar M,\
			\bar g(\xi,\mathrm df_p(v))=0
			\text{ 对所有 }v\in\mathrm{T}_pM\}.
	\]
	其截面称为法向量场。正交分解记为
	\[
		f^*\mathrm{T}\bar M=\mathrm df(\mathrm{T}M)\oplus N_fM,
		\qquad W=W^\top+W^\perp.
	\]
\end{definition}
\symidx{\(N_fM\)}

这个分解随点光滑。事实上，在源坐标中令 \(f_i=\mathrm df(\partial_i)\)，则切向投影为
\[
	W^\top=\sum_{i,j}g^{ij}\bar g(W,f_i)f_j.
\]
矩阵 \((g^{ij})\) 光滑，故两个投影均光滑，法向投影具有常秩 \(r\)。局部选择线性无关的投影像，再作光滑正交化，就得到法丛的局部正交单位标架。

\ProseParagraphBreak

沿 \(f\) 的环境协变导数采用\cref{g3:prop:induced-connection-projection}中的定义。若 \(f^a=y^a\circ f\)、\(W=W^a\partial_{y^a}|_f\)，则
\[
	\bar\nabla^f_XW
	=\left(X(W^a)+\bar\Gamma^a_{bc}(f)\,X(f^b)W^c\right)
		\partial_{y^a}|_f.
\]
它只使用沿 \(f\) 的数据。以下将 \(\bar\nabla^f_X\) 简写为 \(\bar\nabla_X\)，其中 \(X\) 始终是源流形上的向量场。这一联络满足环境度量的乘积规则，而上述命题给出
\[
	(\bar\nabla_X\mathrm df(Y))^\top=\mathrm df(\nabla_XY).
\]

\begin{definition}\label{g3:def:second-fundamental-form}
	设 \(f:(M^n,g)\to(\bar M,\bar g)\) 为光滑等距浸入。定义它的\term[dierjibenxingshi]{第二基本形式}[second fundamental form]为法丛值的二阶协变张量
	\[
		B(X,Y)\coloneqq(\bar\nabla_X\mathrm df(Y))^\perp.
	\]
\end{definition}
\symidx{\(B\)}

\begin{proposition}\label{g3:prop:gauss-decomposition}
	光滑等距浸入 \(f:(M,g)\to(\bar M,\bar g)\) 的第二基本形式光滑、对称，并对两个输入都为 \(C^\infty(M)\) 线性。切向与法向分解给出
	\begin{equation}\label{g3:eq:gauss-decomposition}
		\bar\nabla_X\mathrm df(Y)=\mathrm df(\nabla_XY)+B(X,Y).
	\end{equation}
\end{proposition}
\begin{proof}
	对光滑函数 \(h\)，第一输入的函数线性直接来自联络；第二输入则满足
	\[
		\bar\nabla_X\mathrm df(hY)
		=X(h)\mathrm df(Y)+h\bar\nabla_X\mathrm df(Y).
	\]
	第一项切向，取法向投影后为零，故 \(B(X,hY)=hB(X,Y)\)。各投影和联络系数光滑，因而 \(B\) 光滑。

	环境联络的无挠性沿浸入表现为
	\[
		\bar\nabla_X\mathrm df(Y)-\bar\nabla_Y\mathrm df(X)
		=\mathrm df([X,Y]).
	\]
	可由上面的坐标公式直接核对：含 \(\bar\Gamma^a_{bc}\) 的两项因下指标对称而相消，余下分量为 \(X(Yf^a)-Y(Xf^a)=[X,Y]f^a\)。右侧切向，取法向投影即得 \(B(X,Y)=B(Y,X)\)。最后将环境导数分成两个正交分量，切向部分由诱导联络公式给出，得到式\eqref{g3:eq:gauss-decomposition}。
\end{proof}

式\eqref{g3:eq:gauss-decomposition}称为 Gauss 公式。它把内在联络与浸入的额外二阶数据分开。反过来，对法向场求导，也会出现切向分量；由于切向场与法向场的内积恒为零，这个分量已由 \(B\) 决定。

\begin{definition}\label{g3:def:shape-operator}
	设 \(f:(M,g)\to(\bar M,\bar g)\) 为光滑等距浸入。对法向量 \(\xi\in N_{f,p}M\)，定义\term[xingsuanzi]{形算子}[shape operator]
	\(A_\xi:\mathrm{T}_pM\to\mathrm{T}_pM\)，使
	\begin{equation}\label{g3:eq:shape-operator-pairing}
		g(A_\xi X,Y)=\bar g(B(X,Y),\xi)
		\qquad(X,Y\in\mathrm{T}_pM).
	\end{equation}
	正定内积保证这个线性算子存在且唯一。
\end{definition}
\symidx{\(A_\xi\)}

\begin{proposition}\label{g3:prop:weingarten-decomposition}
	设 \(f:(M,g)\to(\bar M,\bar g)\) 为光滑等距浸入。形算子随 \(\xi\) 线性变化，且每个 \(A_\xi\) 都自伴。对光滑法向场 \(\xi\)，令
	\(\nabla^\perp_X\xi=(\bar\nabla_X\xi)^\perp\)，则
	\begin{equation}\label{g3:eq:weingarten-formula}
		\bar\nabla_X\xi=-\mathrm df(A_\xi X)+\nabla^\perp_X\xi.
	\end{equation}
\end{proposition}
\begin{proof}
	关于 \(\xi\) 的线性由式\eqref{g3:eq:shape-operator-pairing}及唯一性得到。在坐标中，\((A_\xi)^i{}_j=g^{ik}\bar g(B(\partial_j,\partial_k),\xi)\)，故对光滑法向场，形算子场也光滑。由 \(B\) 对称，
	\[
		g(A_\xi X,Y)=\bar g(B(X,Y),\xi)
		=\bar g(B(Y,X),\xi)=g(X,A_\xi Y),
	\]
	所以 \(A_\xi\) 自伴。对恒等式 \(\bar g(\xi,\mathrm df(Y))=0\) 沿 \(X\) 求导，得到
	\[
		\bar g(\bar\nabla_X\xi,\mathrm df(Y))
		=-\bar g(\xi,\bar\nabla_X\mathrm df(Y))
		=-g(A_\xi X,Y).
	\]
	这对所有 \(Y\) 成立，故切向分量是 \(-\mathrm df(A_\xi X)\)；余下分量正是所定义的 \(\nabla^\perp_X\xi\)。
\end{proof}

式\eqref{g3:eq:weingarten-formula}称为\term[weingartengongshi]{Weingarten 公式}[Weingarten formula]。\footnote{Julius Weingarten（1836--1910），德国数学家，研究曲面等距变形与曲率线。}其负号来自正交关系的求导。若余维为一，取局部单位法向 \(\nu\)，则 \(\bar\nabla_X\nu\) 与 \(\nu\) 正交，已没有法向分量，所以 \(\bar\nabla_X\nu=-\mathrm df(A_\nu X)\)。

\begin{example}\label{g3:exm:sphere-extrinsic-sign}
	对半径 \(a>0\) 的圆球面 \(S^n(a)\subset\mathbb R^{n+1}\)，取外向单位法向 \(\nu(p)=p/a\)。环境联络为普通导数，故对切向量 \(X\)，
	\[
		\mathrm D_X\nu=\frac1aX,\qquad
		A_\nu=-\frac1a\operatorname{id},\qquad
		B(X,Y)=-\frac1a g(X,Y)\nu.
	\]
	这些等式直接来自法向的变化，与第一章球面诱导联络的符号相容。把 \(\nu\) 改为 \(-\nu\) 会改变 \(A_\nu\) 的符号，但不改变法丛值张量 \(B\)。
\end{example}

\subsection{法联络及其曲率}
\label{g3:subsec:0401:02}

法向导数中的切向部分已经由 \(B\) 决定，余下部分比较的是相邻法向空间中的向量。余维为一时，单位法向的这部分导数为零；余维较高时，保持长度的法向量还可以在法向空间内部转动。

\begin{definition}\label{g3:def:normal-connection}
	设 \(f:(M,g)\to(\bar M,\bar g)\) 为光滑等距浸入。法丛上的\term[falianluo]{法联络}[normal connection]定义为
	\[
		\nabla^\perp_X\xi=(\bar\nabla_X\xi)^\perp,
		\qquad \xi\in\Gamma(N_fM).
	\]
	法向场 \(\xi\) 若满足 \(\nabla^\perp_X\xi=0\) 对所有 \(X\) 成立，则称为法平行的。
\end{definition}
\symidx{\(\nabla^\perp\)}

\begin{proposition}\label{g3:prop:normal-connection-metric}
	光滑等距浸入的 \(\nabla^\perp\) 是法丛上的光滑联络，并与法丛上的诱导内积相容，即
	\begin{equation}\label{g3:eq:normal-connection-metric}
		X\bar g(\xi,\zeta)
		=\bar g(\nabla^\perp_X\xi,\zeta)
			+\bar g(\xi,\nabla^\perp_X\zeta).
	\end{equation}
\end{proposition}
\begin{proof}
	正交投影为光滑丛映射，环境联络对求导方向为函数线性，所以
	\[
		\nabla^\perp_{hX}\xi=h\nabla^\perp_X\xi,\qquad
		\nabla^\perp_X(h\xi)=X(h)\xi+h\nabla^\perp_X\xi.
	\]
	与实线性一起，这正是丛联络的各项公理。环境度量相容性给出
	\[
		X\bar g(\xi,\zeta)
		=\bar g(\bar\nabla_X\xi,\zeta)+\bar g(\xi,\bar\nabla_X\zeta).
	\]
	这里 \(\xi,\zeta\) 都法向，另一个输入的切向分量不贡献内积，因此可将两个环境导数都替换为其法向投影，得到式\eqref{g3:eq:normal-connection-metric}。
\end{proof}

\begin{definition}\label{g3:def:normal-curvature}
	光滑等距浸入的\term[faqulv]{法曲率}[normal curvature]为法丛联络的曲率
	\begin{equation}\label{g3:eq:normal-curvature-definition}
		R^\perp(X,Y)\xi
		\coloneqq\nabla^\perp_X\nabla^\perp_Y\xi
			-\nabla^\perp_Y\nabla^\perp_X\xi
			-\nabla^\perp_{[X,Y]}\xi.
	\end{equation}
	若 \(R^\perp=0\)，称法联络平坦。
\end{definition}
\symidx{\(R^\perp\)}

\begin{proposition}\label{g3:prop:normal-connection-curvature}
	设 \(f:(M,g)\to(\bar M,\bar g)\) 为光滑等距浸入。\(R^\perp(X,Y)\xi\) 对 \(X,Y,\xi\) 均为函数线性，对 \(X,Y\) 反对称，且在法向输入上满足
	\begin{equation}\label{g3:eq:normal-curvature-skew}
		\bar g(R^\perp(X,Y)\xi,\zeta)
		=-\bar g(\xi,R^\perp(X,Y)\zeta).
	\end{equation}
	因而 \(R^\perp\) 是取值于法丛斜自伴算子的二形式。
\end{proposition}
\begin{proof}
	反对称性来自定义。对光滑函数 \(h\)，利用
	\([hX,Y]=h[X,Y]-Y(h)X\)，展开第一输入得
	\[
		\begin{aligned}
			R^\perp(hX,Y)\xi
			={}&h\nabla^\perp_X\nabla^\perp_Y\xi
				-\nabla^\perp_Y(h\nabla^\perp_X\xi)\\
			&-h\nabla^\perp_{[X,Y]}\xi
				+Y(h)\nabla^\perp_X\xi
			=hR^\perp(X,Y)\xi.
		\end{aligned}
	\]
	第二输入的函数线性由反对称性得到。对于法向输入，
	\[
		\nabla^\perp_X\nabla^\perp_Y(h\xi)
		=X(Yh)\xi+(Yh)\nabla^\perp_X\xi
			+(Xh)\nabla^\perp_Y\xi
			+h\nabla^\perp_X\nabla^\perp_Y\xi.
	\]
	交换 \(X,Y\) 并减去括号方向的导数后，函数的二阶导数合成为
	\(XYh-YXh-[X,Y]h=0\)，两个交叉项也抵消，因此第三输入函数线性。局部标架展开说明它只取决于各输入在一点的值，并且光滑。

	最后将度量相容式连续使用两次：
	\[
		\begin{aligned}
			XY\bar g(\xi,\zeta)
			={}&\bar g(\nabla^\perp_X\nabla^\perp_Y\xi,\zeta)
				+\bar g(\nabla^\perp_Y\xi,\nabla^\perp_X\zeta)\\
			&+\bar g(\nabla^\perp_X\xi,\nabla^\perp_Y\zeta)
				+\bar g(\xi,\nabla^\perp_X\nabla^\perp_Y\zeta).
		\end{aligned}
	\]
	交换 \(X,Y\) 后中间两项相消，再减去 \([X,Y]\bar g(\xi,\zeta)\)，左侧为零，右侧为式\eqref{g3:eq:normal-curvature-skew}两边移到一起的和。
\end{proof}

在局部正交单位法标架 \(\nu=(\nu_1,\ldots,\nu_r)\) 中，定义一形式矩阵 \(\eta=(\eta^\alpha{}_\beta)\)，使
\begin{equation}\label{g3:eq:normal-connection-matrix}
	\nabla^\perp\nu_\beta=\sum_\alpha\nu_\alpha\eta^\alpha{}_\beta,
	\qquad \eta^\alpha{}_\beta(X)=\bar g(\nabla^\perp_X\nu_\beta,\nu_\alpha).
\end{equation}
度量相容性给出 \(\eta^{\mathsf T}=-\eta\)。若法向场的列分量为 \(u\)，即 \(\xi=\nu u\)，则
\[
	\nabla^\perp_X\xi=\nu\bigl(Xu+\eta(X)u\bigr).
\]
因此平行性是关于各方向导数的方程 \(Xu=-\eta(X)u\)。

\begin{proposition}\label{g3:prop:normal-curvature-matrix}
	光滑等距浸入在局部正交单位法标架中的法曲率矩阵为
	\begin{equation}\label{g3:eq:normal-curvature-matrix}
		F^\perp=\mathrm d\eta+\eta\wedge\eta,\qquad
		R^\perp(X,Y)\nu_\beta
			=\sum_\alpha\nu_\alpha(F^\perp)^\alpha{}_\beta(X,Y).
	\end{equation}
	这里 \((\eta\wedge\eta)^\alpha{}_\beta
	=\sum_\gamma\eta^\alpha{}_\gamma\wedge\eta^\gamma{}_\beta\)。
\end{proposition}
\begin{proof}
	将式\eqref{g3:eq:normal-connection-matrix}代入曲率定义，\(\nu_\alpha\) 的系数为
	\[
		\begin{aligned}
			&X(\eta^\alpha{}_\beta(Y))
				-Y(\eta^\alpha{}_\beta(X))
				-\eta^\alpha{}_\beta([X,Y])\\
			&\qquad+\sum_\gamma
				\bigl(\eta^\alpha{}_\gamma(X)\eta^\gamma{}_\beta(Y)
				-\eta^\alpha{}_\gamma(Y)\eta^\gamma{}_\beta(X)\bigr).
		\end{aligned}
	\]
	第一行是 \(\mathrm d\eta^\alpha{}_\beta(X,Y)\)，第二行是矩阵外积的对应分量，故得到所述公式。
\end{proof}

局部正交单位法标架只要求每点的内积矩阵为单位矩阵。若进一步要求它在一个开邻域内法平行，就要求 \(\eta\) 在整个邻域为零，由式\eqref{g3:eq:normal-curvature-matrix}必有 \(R^\perp=0\)。下面的直接计算说明，这一要求对一般浸入不成立。

\begin{example}\label{g3:exm:nonflat-normal-connection}
	在 \(\mathbb R^4\) 中考虑图像浸入
	\[
		f(x,y)=\left(x,y,\frac{x^2}{2},xy\right).
	\]
	其坐标切向量为 \(f_x=(1,0,x,y)\)、\(f_y=(0,1,0,x)\)，以下两个场处处法向且线性无关：
	\[
		\xi_1=(-x,0,1,0),\qquad
		\xi_2=(-y,-x,0,1).
	\]
	它们不要求在整个邻域正交，在原点则为正交单位基。法向 Gram 矩阵及其行列式为
	\[
		G_N=
		\begin{pmatrix}1+x^2&xy\\xy&1+x^2+y^2\end{pmatrix},
		\qquad
		d_N=(1+x^2)(1+x^2+y^2)-x^2y^2>0.
	\]
	由于 \(\partial_y\xi_1=0\)，有 \(\nabla^\perp_{\partial_y}\xi_1=0\)。写
	\(\nabla^\perp_{\partial_x}\xi_1=a\xi_1+b\xi_2\)，对两法向场取内积便得
	\[
		G_N\binom{a}{b}=\binom{x}{y},\qquad
		a=\frac{x(1+x^2)}{d_N},\quad b=\frac{y}{d_N}.
	\]
	在原点 \(a=b=0\)、\(\partial_ya=0\)、\(\partial_yb=1\)，所以
	\[
		R^\perp(\partial_x,\partial_y)\xi_1\big|_{(0,0)}
		=-\nabla^\perp_{\partial_y}(a\xi_1+b\xi_2)\big|_{(0,0)}
		=-\xi_2(0,0)\ne0.
	\]
	因此任何包含原点的邻域都不存在处处法平行的法标架。环境空间本身平坦，并不使这个法联络平坦。
\end{example}

若法丛秩为一，任一局部单位法向 \(\nu\) 满足
\(\bar g(\nabla^\perp_X\nu,\nu)=0\)，而法向空间由 \(\nu\) 张成，故 \(\nabla^\perp\nu=0\)、\(R^\perp=0\)。这个局部结论不需要一个全局单位法向；是否能在整个流形选取这样的法向，是法线丛的另一项整体性质。

\subsection{平均曲率向量}
\label{g3:subsec:0401:03}

第二基本形式为每个切向方向给出一个法向弯曲向量。对正交方向求平均，应当得到不依赖所选切向基的法向量；这正是对 \(B\) 作度量迹的作用。

\begin{definition}\label{g3:def:mean-curvature-vector}
	设 \(f:(M^n,g)\to(\bar M,\bar g)\) 为光滑等距浸入，\(n\ge1\)。其\term[pingjunqulvxiangliang]{平均曲率向量}[mean curvature vector]为
	\begin{equation}\label{g3:eq:mean-curvature-vector}
		H\coloneqq\frac1n\operatorname{tr}_gB
			=\frac1n\sum_{i=1}^nB(e_i,e_i),
	\end{equation}
	其中 \(e_1,\ldots,e_n\) 为任意正交单位基。若 \(H\) 恒为零，称 \(f\) 为\term[jixiaojinru]{极小浸入}[minimal immersion]，相应浸入子流形称为极小子流形。
\end{definition}
\symidx{\(H\)}

若 \(e'_i=\sum_j e_jQ_{ji}\)，其中 \(Q\) 为正交矩阵，则
\[
	\sum_iB(e'_i,e'_i)
	=\sum_{j,k}\left(\sum_iQ_{ji}Q_{ki}\right)B(e_j,e_k)
	=\sum_jB(e_j,e_j).
\]
这既证明定义与正交基无关，也说明 \(H\) 是光滑法向场。在一般坐标中，
\[
	H=\frac1n g^{ij}B(\partial_i,\partial_j).
\]
与任意法向量 \(\xi\) 配对，式\eqref{g3:eq:shape-operator-pairing}又给出
\begin{equation}\label{g3:eq:mean-curvature-shape-trace}
	\bar g(H,\xi)=\frac1n\operatorname{tr}A_\xi.
\end{equation}
因此 \(H=0\) 等价于每个法向形算子都无迹。

\ProseParagraphBreak

对超曲面，选择单位法向 \(\nu\) 后可写 \(H=h\nu\)，其中 \(h=\operatorname{tr}A_\nu/n\) 是带符号的平均曲率。换成 \(-\nu\) 后，\(A_\nu\) 与 \(h\) 同时变号，向量 \(H\) 不变。由\cref{g3:exm:sphere-extrinsic-sign}，外向法向的圆球面满足
\[
	h=-\frac1a,\qquad H=-\frac1a\nu.
\]
平均曲率向量指向球心；它在这里采用平均值归一化，而非未除以维数的迹。

\begin{example}\label{g3:exm:curve-mean-curvature-vector}
	设 \(\gamma:I\to\bar M\) 为单位速光滑曲线，并给 \(I\) 度量 \(\mathrm ds^2\)。记 \(T=\dot\gamma\)。一维域中的单位坐标向量满足 \(\nabla_{\partial_s}\partial_s=0\)，所以 Gauss 公式给出
	\[
		B(\partial_s,\partial_s)=\bar{\mathrm D}_sT,\qquad
		H=\bar{\mathrm D}_sT.
	\]
	单位速条件使右侧与 \(T\) 正交。于是曲线的平均曲率向量就是其曲率向量，而单位速曲线是极小浸入当且仅当它是环境测地线。
\end{example}

\subsection{全测地和全脐子流形}
\label{g3:subsec:0401:04}

两种特殊弯曲情形可以直接用 \(B\) 描述。一种是所有方向都没有法向弯曲；另一种是所有单位切向方向的法向弯曲完全相同。前者对应测地线的保持，后者对应每个法向形算子都是数量算子。

\begin{definition}\label{g3:def:totally-geodesic-umbilical}
	设 \(f:(M^n,g)\to(\bar M,\bar g)\) 为光滑等距浸入，\(n\ge1\)。若 \(B\equiv0\)，称 \(f\) 为\term[quancedijinru]{全测地浸入}[totally geodesic immersion]。若每点 \(p\) 存在法向量 \(\zeta_p\)，使
	\[
		B_p(v,v)=g_p(v,v)\zeta_p\qquad
		\text{对所有 }v\in\mathrm{T}_pM,
	\]
	则称 \(f\) 为\term[quanqijinru]{全脐浸入}[totally umbilical immersion]。
\end{definition}

\begin{proposition}\label{g3:prop:totally-geodesic-characterization}
	设 \(f:(M,g)\to(\bar M,\bar g)\) 为光滑等距浸入。则 \(f\) 全测地，当且仅当每条 \(M\) 中的测地线经 \(f\) 映射后仍是保持原参数的环境测地线。此时，对每个 \(p\) 及足够小的 \(v\in\mathrm{T}_pM\)，有
	\begin{equation}\label{g3:eq:totally-geodesic-exponential}
		f(\exp_p^M v)=\exp_{f(p)}^{\bar M}(\mathrm df_pv).
	\end{equation}
\end{proposition}
\begin{proof}
	对任意光滑曲线 \(\gamma\)，将 Gauss 公式沿曲线使用，得到加速度分解
	\begin{equation}\label{g3:eq:immersion-curve-acceleration}
		\bar{\mathrm D}_t(\mathrm df(\dot\gamma))
		=\mathrm df(\mathrm D_t\dot\gamma)+B(\dot\gamma,\dot\gamma).
	\end{equation}
	具体地，在坐标中写 \(\dot\gamma=\dot x^i\partial_i\)，左侧为
	\(\ddot x^i f_i+\dot x^i\dot x^j\bar\nabla_{\partial_i}f_j\)；
	代入式\eqref{g3:eq:gauss-decomposition}，切向部分合成为
	\(\mathrm df((\ddot x^k+\Gamma^k_{ij}\dot x^i\dot x^j)\partial_k)\)，
	法向部分即上式所列的 \(B\)。

	若 \(B=0\)，当 \(\gamma\) 为测地线时，两项都为零，故 \(f\circ\gamma\) 是环境测地线。反过来，固定 \(p\) 及任意 \(v\in\mathrm{T}_pM\)。由\cref{g3:thm:geodesic-ivp}，存在以 \((p,v)\) 为初值的局部测地线。若其像也是环境测地线，在时间零代入加速度分解，得到 \(B_p(v,v)=0\)。由 \(B\) 对称，
	\[
		2B_p(v,w)=B_p(v+w,v+w)-B_p(v,v)-B_p(w,w)=0.
	\]
	所以 \(B=0\)。

	最后，在全测地情形，\(f\circ\gamma_{p,v}\) 与环境中以
	\((f(p),\mathrm df_pv)\) 为初值的测地线相同，这由同一初值唯一性得到。当 \(v\) 足够小，使两条曲线均定义到时间一时，比较时间一的值即得式\eqref{g3:eq:totally-geodesic-exponential}。
\end{proof}

\begin{proposition}\label{g3:prop:umbilical-characterization}
	设 \(f:(M^n,g)\to(\bar M,\bar g)\) 为光滑等距浸入，\(n\ge1\)。它全脐，当且仅当
	\begin{equation}\label{g3:eq:umbilical-second-fundamental-form}
		B(X,Y)=g(X,Y)H.
	\end{equation}
	这又等价于对每个法向量 \(\xi\)，有
	\[
		A_\xi=\bar g(H,\xi)\operatorname{id}.
	\]
	全脐浸入是全测地的，当且仅当其平均曲率向量为零。
\end{proposition}
\begin{proof}
	若 \(B_p(v,v)=g_p(v,v)\zeta_p\)，对 \(v+w\) 作同样计算，再减去 \(v,w\) 两式，得到
	\(B_p(v,w)=g_p(v,w)\zeta_p\)。在正交单位基中取迹，便有
	\(nH_p=n\zeta_p\)，所以 \(\zeta_p=H_p\)，得到式\eqref{g3:eq:umbilical-second-fundamental-form}。反向蕴含由令 \(Y=X\) 得到。

	将该式与任意 \(\xi\) 配对，形算子的定义给出
	\[
		g(A_\xi X,Y)=g(X,Y)\bar g(H,\xi),
	\]
	由度量非退化性得到数量算子表达。反过来，若所有 \(A_\xi\) 都具有这一形式，则 \(B(X,Y)-g(X,Y)H\) 与每个法向量正交；它本身法向，故为零。最后在全脐公式中令 \(H=0\) 即得 \(B=0\)，而 \(B=0\) 总使其迹 \(H\) 为零。
\end{proof}

Euclid 空间中的仿射子空间有常值切空间，环境导数的法向分量恒为零，因此全测地。圆球面由\cref{g3:exm:sphere-extrinsic-sign}满足 \(B=gH\)，故全脐，但它的 \(H\ne0\)，所以并非全测地。维数一时，每个单位切向量只差一个符号，任何曲线浸入都自动全脐；此时全测地与极小的条件则都要求曲率向量为零。
% 纲要条目排为编号小节标题，并进入目录与书签；说明段仍保留于纲要注释中。
